
\subsection{连续扰动族}
\label{subsec:8-10-3}
\liftIndex{I360}

最后讨论第~\ref{sec:7-5} 节为定义 $A_K$ 同态
\[
 \mathfrak f_\ell\colon
 B_\ell(\mathbb R\mathcal X_L)\longrightarrow\Omega(L)
\]
所用模空间的定向。在第~\ref{sec:7-5} 节中，我们定义了空间
\[
 (L\times\mathcal N'_{\ell+1})
 \mathbin{\times_{L^{\ell+1}}}
 (P'_1\times\cdots\times P'_\ell\times L)
\]
的仓西邻域
\[
 U_{\ell+1}(\mathfrak f,\vec P')
 = (U_{\ell+1}\times\mathcal N'_{\ell+1})
 \mathbin{\times_{L^{\ell+1}}}
 (P'_1\times\cdots\times P'_\ell\times L).
\]
（情形 $U_{k+1}(\mathfrak m,\vec P)$ 的问题相同，故从略。）我们用它
依式~\eqref{eq:7-5-48} 定义 $\mathfrak X_\ell(\vec P')$：
\begin{equation}
 \mathfrak X_\ell(\vec P')
 =\bigl\{(w,p)\in W_\ell\times
 U_{\ell+1}(\mathfrak f,\vec P')
 \bigm|
 \mathfrak t_{\mathfrak f,\ell+1}^{\vec P'}(w,p)=0
 \bigr\}.
 \tag{7.5.48}
\end{equation}
这里 $\mathfrak t_{\mathfrak f,\ell+1}^{\vec P'}$ 是
$W_\ell\times U_{\ell+1}(\mathfrak f,\vec P')$ 上障碍丛的一个截面。
我们有投影
\[
 \pi_{W_\ell}\colon\mathfrak X_\ell(\vec P')\longrightarrow W_\ell.
\]

我们在 $W_\ell$ 上取了一个顶次形式 $\omega_\ell$，并定义
\begin{equation}
 \mathfrak f_\ell(\vec P')
 =ev_!\bigl(\pi_{W_\ell}^*(\omega_\ell)\bigr)\in\Omega(L).
 \tag{7.5.49}
\end{equation}
这里 $ev\colon\mathfrak X_\ell(\vec P')\to L$ 是第 $0$ 个标记点处的
评价映射。

$W_\ell$ 的定向由 $\omega_\ell$ 给出。我们需要找出
$U_{\ell+1}(\mathfrak f,\vec P')$ 的一个定向，使它诱导
$\mathfrak X_\ell(\vec P')$ 的定向，从而确定上式~\eqref{eq:7-5-49}
的符号。我们将把这个问题约化为定义滤过 $A_\infty$ 同态所用模空间上的
定向问题；后一个问题已在第~\ref{subsec:8-9-1} 小节讨论。

令 $d\lvert\omega_\ell\rvert$ 为 $W_\ell$ 上满足
\[
 \int_{W_\ell}f\,\omega_\ell
 =\int_{W_\ell}f\,d\lvert\omega_\ell\rvert
\]
的光滑测度，其中 $f$ 是任意正函数。由标准横截性定理，可以找到
$W_\ell$ 中一个具有全 $d\lvert\omega_\ell\rvert$ 测度的子集
$W_\ell^0$，使得对每个 $w\in W_\ell^0$，子集
\[
 \pi_{W_\ell}^{-1}(w)\subset\mathfrak X_\ell(\vec P')
\]
都是光滑子流形。若 $\pi_{W_\ell}^{-1}(w)$ 已定向，则
$ev(\pi_{W_\ell}^{-1}(w))\subset L$ 定义 $L$ 上的一个流。而且不难证明
\begin{equation}
 ev_!\bigl(\pi_{W_\ell}^*(\omega_\ell)\bigr)
 =\int_{w\in W_\ell^0}
 ev\bigl(\pi_{W_\ell}^{-1}(w)\bigr)
 \,d\lvert\omega_\ell\rvert(w)
 \tag{8.10.8}\label{eq:8-10-8}
\end{equation}
除一个符号外，作为光滑微分形式的等式成立。因此，要为
式~\eqref{eq:7-5-49} 找到适当的符号，便约化为给
$\pi_{W_\ell}^{-1}(w)$ 找到适当的定向。

下面回忆，$W_\ell$ 是参数化
$U_{\ell+1}(\mathfrak f,\vec P')$ 上仓西结构之扰动（截面）族的空间。
集合
\[
 \pi_{W_\ell}^{-1}(w)
 =\mathfrak X_\ell(\vec P')\cap
 \bigl(\{w\}\times U_{\ell+1}(\mathfrak f,\vec P')\bigr)
\]
是截面
\[
 p\longmapsto
 \mathfrak t_{\mathfrak f,\ell+1,w}^{\vec P'}(p)
 :=\mathfrak t_{\mathfrak f,\ell+1}^{\vec P'}(w,p)
\]
的零点集。若把 $\mathfrak t_{\mathfrak f,\ell+1}^{\vec P'}$ 看作
$U_{\ell+1}(\mathfrak f,\vec P')$ 上由 $W_\ell$ 参数化的截面族
\[
 \bigl\{\mathfrak t_{\mathfrak f,\ell+1,w}^{\vec P'}
 \bigr\}_{w\in W_\ell},
\]
那么 $\pi_{W_\ell}^{-1}(w)$ 就是其中相应截面的零点集。

考虑空间
\[
 U_{\ell+1}(\mathfrak f,\vec P')^+
 =(U_{\ell+1}\times\mathcal N_{\ell+1})
 \mathbin{\times_{L^{\ell+1}}}
 (P'_1\times\cdots\times P'_\ell\times L)
\]
以及障碍丛
\[
 E_{\ell+1}(\mathfrak f)^+
 =\mathcal N_{\ell+1}\times E_{\ell+1}.
\]
（注意，在 $U_{\ell+1}(\mathfrak f,\vec P')^+$ 与
$E_{\ell+1}(\mathfrak f)^+$ 的右端，记号 $\mathcal N_{\ell+1}$ 不带撇号。）
它们在空间
\begin{equation}
 (L\times\mathcal N_{\ell+1})
 \mathbin{\times_{L^{\ell+1}}}
 (P'_1\times\cdots\times P'_\ell\times L)
 \tag{8.10.9}\label{eq:8-10-9}
\end{equation}
上定义一个仓西结构。注意，式~\eqref{eq:8-10-9} 的空间连同其仓西结构
可以与下列空间相识别：
\begin{equation}
 \mathcal M_{\ell+1}^{\mathrm{main}}
 \bigl(\{J_\rho\}_\rho:\beta_0;
 \operatorname{top}(\rho);P'_1,\ldots,P'_\ell\bigr).
 \tag{8.10.10}\label{eq:8-10-10}
\end{equation}
我们已经在第~\ref{subsec:8-9-1} 小节描述了
式~\eqref{eq:8-10-10} 上的定向：我们在那里用该空间构造了
$A_\infty$ 同态
\[
 \mathfrak f\colon
 (\mathbb R\mathcal X_L,\overline{\mathfrak m})
 \longrightarrow
 (\mathbb R\mathcal X'_L,\overline{\mathfrak m}'),
\]
其两端是由两种不同的扰动选择和两种不同的奇异单形集
$\mathcal X_L,\mathcal X'_L$ 得到的两个 $A_\infty$ 代数；这些代数由
定理~\ref{thm:3-4-8} 建立。因此，我们已经在第~\ref{sec:8-9} 节给
$U_{\ell+1}(\mathfrak f,\vec P')^+$ 定义了相容的定向。

另一方面，由定理~\ref{thm:7-5-63}，存在一个光滑满射
\[
 \mathcal N_{\ell+1}\longrightarrow\mathcal N'_{\ell+1},
\]
它在边界之外是微分同胚。因此，$U_{\ell+1}(\mathfrak f,\vec P')^+$
的定向诱导 $U_{\ell+1}(\mathfrak f,\vec P')$ 的定向。再按照
约定~\ref{conv:8-2-1} 中式~\eqref{eq:8-2-2} 的规则，我们可以给
$\mathfrak X_\ell(\vec P')$ 定向；这一取向确定了式~\eqref{eq:7-5-49}
的符号。
